\section{Rangkuman}
\label{sec:rangkuman-bab5}

Pada bab ini telah dibahas:

\begin{itemize}
    \item Register adalah tempat penyimpanan tercepat di dalam
          prosesor. Prosesor 8086 memiliki empat belas register
          16 bit yang dapat dipakai programmer.

    \item Register umum \reg{AX}, \reg{BX}, \reg{CX}, dan \reg{DX}
          dapat diakses sebagai dua byte terpisah: \reg{AH}/\reg{AL},
          \reg{BH}/\reg{BL}, \reg{CH}/\reg{CL}, dan \reg{DH}/\reg{DL}.
          Menulis ke byte rendah tidak mengubah byte tinggi.

    \item Setiap register umum memiliki peran khas sebagai akibat
          cara penyandian instruksi: \reg{AX} sebagai akumulator,
          \reg{BX} sebagai basis pengalamatan, \reg{CX} sebagai
          pencacah, dan \reg{DX} sebagai pasangan data pada operasi
          32 bit serta penunjuk alamat port.

    \item Register segmen \reg{CS}, \reg{DS}, \reg{SS}, dan \reg{ES}
          menyimpan nomor segmen, bukan data. Pasangannya adalah
          \reg{IP} untuk kode, \reg{BP}/\reg{SP} untuk \eng{stack},
          serta \reg{SI} dan \reg{DI} untuk data dan tujuan operasi
          string.

    \item Register \eng{pointer} dan \eng{index} menyimpan
          \eng{offset}. \reg{SP} menunjuk puncak \eng{stack} dan
          berubah otomatis pada \asm{PUSH}/\asm{POP}, \reg{BP}
          menjadi acuan tetap di dalam \eng{stack}, sedangkan
          \reg{SI} dan \reg{DI} menunjuk sumber dan tujuan operasi
          string.

    \item Register \reg{IP} menyimpan \eng{offset} instruksi
          berikutnya dan selalu berpasangan dengan \reg{CS}.

    \item Register \reg{FLAGS} 16 bit memuat sembilan \eng{flag}
          aktif: \flag{CF}, \flag{PF}, \flag{AF}, \flag{ZF},
          \flag{SF}, \flag{OF} sebagai \eng{flag} status, serta
          \flag{TF}, \flag{IF}, \flag{DF} sebagai \eng{flag}
          kendali.

    \item \flag{CF} melaporkan simpanan keluar pada bilangan tak
          bertanda, sedangkan \flag{OF} melaporkan \eng{overflow}
          pada bilangan bertanda. Keduanya harus dibedakan dengan
          cermat.

    \item \eng{Flag} menjadi berguna ketika dibaca oleh lompatan
          bersyarat, sedangkan \asm{MOV}, \asm{LEA}, \asm{PUSH}, dan
          \asm{POP} tidak mengubah \eng{flag}.

    \item Register \reg{CS}, \reg{IP}, \reg{SS}, dan \reg{SP}
          mengendalikan alur program dan letak \eng{stack}, sehingga
          tidak boleh diubah sembarangan.
\end{itemize}

\begin{checklistbox}
Tandai setiap butir yang sudah Anda kuasai.

\begin{itemize}
    \item[\checkbox] Saya dapat menyebutkan keempat belas register
          8086 beserta kelompoknya.
    \item[\checkbox] Saya dapat menjelaskan akses byte rendah dan
          byte tinggi pada \reg{AX}, \reg{BX}, \reg{CX}, \reg{DX}.
    \item[\checkbox] Saya dapat menjelaskan peran khas setiap
          register umum beserta alasannya.
    \item[\checkbox] Saya dapat menjelaskan peran \reg{CS}, \reg{DS},
          \reg{SS}, dan \reg{ES} dalam pengalamatan memori.
    \item[\checkbox] Saya dapat menjelaskan perbedaan tugas \reg{SP},
          \reg{BP}, \reg{SI}, \reg{DI}, dan \reg{IP}.
    \item[\checkbox] Saya dapat menyebutkan sembilan \eng{flag}
          beserta posisi bitnya.
    \item[\checkbox] Saya dapat memprediksi \flag{CF}, \flag{OF},
          \flag{SF}, dan \flag{ZF} pada operasi \asm{ADD}.
    \item[\checkbox] Saya dapat menjelaskan akibat mengubah
          \reg{CS}, \reg{IP}, \reg{SS}, atau \reg{SP} secara
          sembarangan.
\end{itemize}

\vspace{0.4em}
Bila masih ada butir yang belum bertanda, ulangi bagian yang
bersesuaian sebelum melanjutkan ke bab berikutnya.
\end{checklistbox}

\begin{infobox}
\textbf{Rujukan bab.} Uraian mengenai fungsi dan penyandian register
8086 bersumber dari \cite{intel1979}. Penjelasan tentang \eng{flag}
dan pengaruhnya terhadap eksekusi diperkuat oleh \cite{brey2009}
dan \cite{hyde2010}. Ringkasan praktis mengenai daftar register dan
\eng{flag} dapat ditelusuri pada \cite{wikibooksx86}.
\end{infobox}
